\section{Read and write management}

As a requirement of the project, reads must be locally performed, while writes have to be approved by the majority, and only then the coordinator can proceed to apply and broadcast the change. The database is a map of \texttt{<idx: Integer, value: Integer>}, and each replica affects only its own database.

\subsection{Assumptions}

Some assumptions must be considered in the following protocol. Some of them are required by the assignment, others were defined by us during the implementation:

\begin{itemize}
\item reliable FIFO links are fundamental assumptions for the system network;

\item the quorum is always available, meaning that a write request is always processed by the coordinator, if the latter does not crash in the meanwhile;

\item the coordinator processes one write at a time, hence a replica can be at most one write behind the coordinator;

\item Akka is always used to notify all the replicas, including the sender itself. However, in this scenario, the plain Akka \texttt{.tell()} function is adopted avoiding latency (that internally the same replica is not plausible), instead of the \texttt{tell()} version in \texttt{AbstractReplica} that leverages the custom network system with delays. One example is the write request forwarding to the coordinator;

    \begin{lstlisting}
if (this.id == this.coordinatorId) {
  // If the request is done to the coordinator itself,
  // send the request using Akka .tell() (without delay)
  getSelf().tell(writeRequestToCoordinator, getSelf());
} else {
  tell(writeRequestToCoordinator, this.replicasGroup.get(this.coordinatorId));
}
    \end{lstlisting}

\end{itemize}

\subsection{Read request}

A client may ask a replica to read a value in the database associated with an index. Hence, the client sends the request and sets a timeout in order to detect the replica failure. If no result is provided before the timeout expiration, the replica is considered crashed and the relative callback is invoked. As soon as the replica receives the request, it immediately reads the database and sends back the result.

\subsection{Write request}

A client may ask a replica to write a value into the database, providing the index and the new value. Similarly to the read request, the client sets a timeout to detect the replica crash. Note that this timeout has to be larger because of the non-local procedure, unlike reads. The timeout, indeed, is significantly affected by the estimated latency (in our project \texttt{AbstractReplica.MAX\_LATENCY} is adopted) and the number of initial replicas.

Each replica stores a queue of \texttt{<client: ActorRef, request: WriteRequest>} pairs. These are the lists of unstable requests that are waiting to be processed by the coordinator and sent back to the respective client. After pushing the request into the queue, a \texttt{WriteRequestToCoordinator} message is sent to the leader. The latter has a queue of \texttt{<request: WriteRequesToCoordinator>} which is the ordered list of unstable writes. As a request arrives, the coordinator pushes it into the queue and always processes the first one by the \texttt{.poll()} function. This operation is important to preserve \textbf{sequential consistency}.

Writes are processed sequentially: a write is taken in charge, a \texttt{WriteNotification} message is broadcasted waiting for the relative acknowledgments, the write is stabilized after receiving ACKs from the majority, and finally the \texttt{WriteOK} message is sent to all replicas such that all of them can perform the write. Moreover, the replica that receives a \texttt{WriteOK} for a request in its queue, sends the result back to the relative client and pops out the element from the queue. The most relevant aspects of this protocol are described in the following sections.

\subsubsection{Write notification acknowledgments}

After a write notification broadcast, a timeout is set for every replica in order to detect eventual failures. A \texttt{pendingAcks<replicaId: Integer>} list is populated with the IDs of all the replicas but the coordinator. When an ACK is received, the ID of the corresponding replica is removed from the list, and when the timeout expires the replicas with the IDs remaining in \texttt{pendingAcks} are considered crashed (no ACK received). Finally, each timeout is associated with an incremental counter called \texttt{acksRound}. This is useful because, in case the coordinator is able to rapidly process two writes, stale ACKs are not considered (their round is smaller than the current one).

\subsubsection{Replicas failures and quorum update}

The quorum is defined as $Q = \lceil N / 2 \rceil + 1$ where $N$ is the number of active replicas. It needs to be updated when a replica crashes thus, during a write notification, the coordinator waits for all ACKs before processing a new write even though it reached the quorum. If some timeouts expire, the relative replicas are considered crashed and the quorum is updated. These replicas are also removed from the \texttt{replicasGroup} which is a local clone of the immutable map passed in the \texttt{initSystem()} method.

\subsubsection{Current write process}

When the current write is completed, the coordinator calls the \texttt{processNextWriteIfAny()} function, which executes the protocol for the first element in the \texttt{writeRequestsQueue} if present. This method is invoked in the following two situations:

\begin{itemize}
\item all the acknowledgments are received, \texttt{WriteOK} was already sent to all replicas and so another write can be processed;

\item the timeout has expired and some ACKs were not received. The \texttt{replicasGroup} and the quorum are updated, and another write can be performed.
\end{itemize}

\subsection{Timeouts setup}

As soon as a client is asked to request a read or a write, it immediately sends the relative message to the replica. A timeout is set for each request in order to detect a replica failure and so two counters were added: \texttt{readRequestCount} increases at every read result and decreases at every read timeout expiration; if a timeout expires and the counter is negative, this means a result did not arrive and the replica crashed. \texttt{writeRequestCount} behaves similarly for write requests.

Both on the client side and the replica side, the following timeouts have been set:

\begin{itemize}
\item the client sends a request and sets a different timeout according to the type of the expected result. The \texttt{readTimeout} depends more on latency, while \texttt{writeTimeout} is more affected by the initial number of replicas in the system and an estimated time needed to process a write;

\item the coordinator sets a timeout for the acknowledgment of the write notification, which is a combination of the latency (plus tolerance) and the number of active actors in the system.
\end{itemize}
